\chapter{Allied Health}

\section*{Definition and Overview}
Allied health is the collective term for a diverse group of health professionals who are not doctors, nurses, or dentists but who are essential to modern healthcare. It includes physiotherapists, occupational therapists, dietitians, speech pathologists, psychologists, podiatrists, exercise physiologists, radiographers, and many others. These professionals contribute to prevention, assessment, treatment, rehabilitation, and the long-term management of chronic disease. Allied health care works alongside, and often complements, medical and nursing care; it is frequently the key to recovering from injury or surgery, living well with a long-term condition, and maintaining independence in daily life. Services are delivered in hospitals, community clinics, private practices, schools, workplaces, and aged care facilities.

\section*{Epidemiology}
Almost everyone will need some form of allied health care during their lifetime. Demand is growing steadily as populations age, rates of chronic disease rise, and the value of rehabilitation and preventive care becomes better recognised. Physiotherapy, podiatry, and occupational therapy are commonly used after injury or surgery and for conditions such as arthritis, diabetes, stroke, and heart disease. Speech pathology is frequently required after stroke and for developmental communication problems, while psychology and dietetics services are increasingly in demand for mental health and lifestyle-related conditions. Access to allied health services varies considerably by region and income, which remains an important equity issue in many health systems.

\section*{Aetiology and Pathophysiology}
The need for allied health care arises from a wide range of situations:

\begin{itemize}
    \item \textbf{Injury and surgery:} restoring movement, strength, and function after fractures, joint replacement, and soft tissue injury.
    \item \textbf{Chronic disease:} diabetes, heart disease, chronic lung disease, and arthritis all benefit from structured exercise, dietary management, and self-management support.
    \item \textbf{Neurological conditions:} stroke, Parkinson's disease, and multiple sclerosis require physiotherapy, occupational therapy, and speech pathology for rehabilitation.
    \item \textbf{Disability and developmental delay:} children with motor, speech, or learning difficulties may need occupational therapy, speech pathology, and physiotherapy.
    \item \textbf{Mental health:} psychologists and counsellors provide treatment for anxiety, depression, and stress.
    \item \textbf{Ageing:} falls prevention, mobility support, and help with daily activities preserve independence.
    \item \textbf{Communication and swallowing:} speech pathology addresses speech, language, voice, and swallowing disorders of many causes.
\end{itemize}

\section*{Clinical Features}
Signs that allied health care may help include:

\begin{itemize}
    \item \textbf{Movement:} persistent pain, weakness, stiffness, or difficulty moving after injury or illness.
    \item \textbf{Daily function:} trouble with dressing, cooking, bathing, or household tasks.
    \item \textbf{Communication:} difficulty speaking, understanding, reading, or writing.
    \item \textbf{Swallowing:} coughing or choking with food or drink, or avoiding certain foods because of swallowing difficulty.
    \item \textbf{Nutrition:} weight problems, dietary concerns, or food allergies needing professional guidance.
    \item \textbf{Feet:} foot pain or foot problems, especially in people with diabetes.
    \item \textbf{Mental wellbeing:} anxiety, low mood, or stress affecting daily life.
    \item \textbf{Balance and falls:} poor balance, recent falls, or fear of falling in older age.
\end{itemize}

\section*{Differential Diagnosis}
Different allied health professionals address different problems, and the appropriate service depends on the specific need:

\begin{itemize}
    \item \textbf{Physiotherapy:} muscle, joint, spinal, and movement problems.
    \item \textbf{Occupational therapy:} practical skills, everyday function, and adaptation after injury or illness.
    \item \textbf{Dietetics:} nutrition, food allergies, eating disorders, and diet-related conditions such as diabetes.
    \item \textbf{Speech pathology:} communication, language, voice, and swallowing disorders.
    \item \textbf{Podiatry:} disorders of the feet, ankles, and lower legs.
    \item \textbf{Psychology:} mental health, behaviour, and psychological support.
    \item \textbf{Exercise physiology:} structured exercise prescription for chronic and complex conditions.
\end{itemize}

A medical review first is often useful to confirm the diagnosis and direct the person to the right professional.

\section*{When to Seek Medical Care}
Allied health services are usually accessed through a doctor, who can confirm the diagnosis and provide a referral or treatment plan. Seek help early if pain, movement, daily function, speech, swallowing, diet, or mood problems are affecting your life, because early intervention generally produces better outcomes.

\textbf{Call 000 (or local emergency services) for:}
\begin{itemize}
    \item Sudden difficulty speaking, weakness of one side, or a drooping face (possible stroke).
    \item Sudden inability to swallow or difficulty breathing.
    \item Chest pain, severe pain, or deformity after an injury.
    \item Collapse or loss of consciousness.
\end{itemize}

\section*{Diagnosis and Investigations}
Allied health assessment is thorough and practical, focusing on function as well as diagnosis:

\begin{itemize}
    \item \textbf{Interview:} a detailed history of the problem, the person's goals, and their home, work, and social situation.
    \item \textbf{Physical assessment:} evaluation of strength, range of motion, balance, gait, or swallowing, as appropriate to the discipline.
    \item \textbf{Functional testing:} measuring performance of everyday tasks, such as walking distance, stair climbing, or meal preparation.
    \item \textbf{Specialised tools:} for example, gait analysis and foot assessment in podiatry, or standardised language testing in speech pathology.
    \item \textbf{Outcome measures:} validated scales that track progress over the course of treatment.
    \item \textbf{Liaison:} sharing findings and plans with the general practitioner and the wider care team.
\end{itemize}

\section*{Management}
Allied health treatment is goal-based, practical, and often delivered in a course of sessions:

\begin{itemize}
    \item \textbf{Exercise programs:} prescribed physiotherapy or exercise physiology to rebuild strength, mobility, and fitness.
    \item \textbf{Hands-on treatment:} joint mobilisation, soft tissue massage, and manual therapy where indicated.
    \item \textbf{Functional retraining:} occupational therapy to restore everyday skills and to select equipment and aids.
    \item \textbf{Nutrition advice:} dietetic plans for weight management, diabetes, allergies, and disease prevention.
    \item \textbf{Communication therapy:} speech pathology for speech, language, voice, or swallowing.
    \item \textbf{Talking therapies:} psychological treatment for mental health conditions and adjustment to illness.
    \item \textbf{Home and equipment advice:} aids, splints, orthotics, and home modifications to improve safety and independence.
    \item \textbf{Review and discharge:} progress checks, outcome measurement, and a written plan for maintaining gains.
\end{itemize}

\section*{Complications}
Delaying or avoiding allied health care can have significant consequences:

\begin{itemize}
    \item Slower and less complete recovery after injury or surgery.
    \item Progressive loss of strength, mobility, and independence.
    \item Falls and fractures in older people.
    \item Worsening control of chronic conditions such as diabetes and lung disease.
    \item Longer hospital stays and greater difficulty returning home after illness.
    \item Ongoing pain or disability that could have improved with earlier treatment.
\end{itemize}

\section*{Prognosis and Outcomes}
With appropriate allied health care, most people recover well from injury, surgery, and illness, and many improve substantially even with long-standing conditions. Outcomes depend on the nature of the condition, how early treatment begins, and how consistently the person follows the program. Chronic conditions respond best to ongoing self-management support, which builds skills that last a lifetime. Accessing the right service at the right time frequently prevents complications, reduces disability, and improves quality of life.

\section*{Prevention and Self-Care}
Keeping active and healthy reduces the need for treatment:

\begin{itemize}
    \item Stay physically active within your ability, and seek advice before starting a new program.
    \item Maintain a balanced diet and a healthy weight.
    \item Practise falls prevention with good footwear, home safety measures, and balance exercises.
    \item Seek early help for persistent pain, foot, speech, or mental health problems.
    \item Follow your allied health professional's home program between sessions.
    \item Ask your doctor about referral options and any entitlements, subsidies, or insurance cover available.
\end{itemize}

\section*{Sources and Further Reading}
The following reputable sources provide reliable information on allied health services and careers:

\begin{itemize}
    \item NHS. \textit{Allied health professions information.} https://www.nhs.uk/
    \item World Health Organization. \textit{Information on rehabilitation and allied health services.} https://www.who.int/
    \item Mayo Clinic. \textit{Health services and physical therapy information.} https://www.mayoclinic.org/
    \item Centers for Disease Control and Prevention. \textit{Physical activity and chronic disease information.} https://www.cdc.gov/
    \item MedlinePlus. \textit{Rehabilitation and therapy health information.} https://medlineplus.gov/
    \item NICE. \textit{Guidance on rehabilitation and care pathways.} https://www.nice.org.uk/
\end{itemize}
